\begin{lemma}[Normalisation; Scheide Lemma 3.11]
Let $(T,e_0,\varphi)$ be globally maximum, with the first two vertices
$x,y$.  There are colors $\alpha\in\overline\varphi(x)$ and
$\beta\in\overline\varphi(y)$ and a normal Tashkinov tree $T'$ on the
same vertex set such that
\[
 h(T')=|V(P_x(\alpha,\beta,\varphi))|.
\]
In the incidence encoding, if the full path has length $L$, this equality
is $h(T')=L/2+1$.
\end{lemma}
\begin{proof}
Write
\[
 x=T.\operatorname{vertex}(0),\qquad
 y=T.\operatorname{vertex}(1).
\]
Unpacking \noderef{GlobalMaximumTriple} gives three facts: $(G,e_0,
\varphi)$ is a $k$-triple, $T$ is a Tashkinov tree for this pair, and
$|V(T)|=t(G)$.  By \noderef{TreeSequence}, the first edge identity of
$T$ is $e_0$, its ends are $x,y$, and $x\ne y$ because the vertex map of
$T$ is injective.  The $k$-triple definition \noderef{KTriple} says in
particular that $\chi'(G)=k+1$ and that $\varphi$ leaves exactly $e_0$
uncolored.  Thus \noderef{IsOneEdgeUncolored} gives
\[
 \varphi(e_0)=\mathsf{none},\qquad
 f\ne e_0\Longrightarrow \varphi(f)=\mathsf{some}(c)
 \text{ for some }c.
\tag{1}
\]

We next verify the fixed-pair maximality needed below.  Let $R$ be a
Tashkinov tree for $(e_0,\varphi)$ with $V(T)\subseteq V(R)$.  Since
$\chi'(G)=k+1$, the integer used in \noderef{TashkinovOrder} is
$\chi'(G)-1=k$.  Therefore $|V(R)|$ is one of the cardinalities over
which $t(G)$ is maximized, and $|V(R)|\le t(G)=|V(T)|$.  The inclusion
$V(T)\subseteq V(R)$ between the finite sets of
\noderef{TreeSequenceVertexSet}, together with this cardinal inequality,
forces equality of the two sets and hence $V(R)\subseteq V(T)$.  This is
every clause of \noderef{FixedPairMaximalTashkinovTree}: the required
Tashkinov-tree clause for $T$ came from global maximality, and the preceding
argument is its maximality clause.

Since $x$ and $y$ are ends of $e_0$, the incident case of
\noderef{ScheideGlobalMaximumMissingCounts} gives at least two missing
colors at each.  Choose
\[
 \alpha\in\overline\varphi(x),\qquad
 \beta\in\overline\varphi(y),
\tag{2}
\]
where missing-color membership has the meaning in
\noderef{MissingColors}.  Apply \noderef{ScheideMaximalKempePath} with
$S=T$, indices $0<1$, the $k$-triple above, the just-proved fixed-pair
maximality, the Tashkinov certificate for $T$, and (2).  It yields
$\alpha\ne\beta$ and a simple incidence path $P$ from $x$ to $y$ such
that
\[
 \{z:z\in P.\operatorname{support}\}
 =C_{\alpha,\beta}(x),\qquad 4\mid L:=|P|,
\tag{3}
\]
and every vertex-node on $P$ belongs to $V(T)$.  Here the component in
(3) is exactly \noderef{TwoColorComponent} in the incidence graph
\noderef{TwoColorIncidenceGraph}.

The graph in \noderef{TwoColorIncidenceGraph} is bipartite between vertex
identities and edge identities, so the successive nodes of $P$, whose two
ends are vertex-nodes, alternate between the two sorts.  Put
\[
 \ell=L/2,\qquad p=\ell+1.
\]
Because $x\ne y$, the path has positive length; together with $4\mid L$
this gives $L\ge4$, $\ell\ge2$, and $p\ge3$.  Reverse the ordered list of
nodes of $P$ and name its identities as
\[
 P^{\mathrm{rev}}=(y=v_1,f_2,v_2,f_3,\ldots,
                    f_{p-1},v_{p-1},f_p,v_p=x).
\tag{4}
\]
Thus (4) has precisely the $p$ vertex identities
$v_1,\ldots,v_p$ and the $p-1=\ell$ edge identities
$f_2,\ldots,f_p$.  Since $P$ is a path, its support has no repeated node;
consequently all the displayed vertex identities are distinct and all the
displayed edge identities are distinct.  For each $2\le r\le p$, the two
incidence adjacencies surrounding $f_r$ in (4), together with
\noderef{TwoColorIncidenceGraph}, say
\[
 \varphi(f_r)=\mathsf{some}(\alpha)
 \text{ or }\mathsf{some}(\beta).
\tag{5a}
\]
They also say that both $v_{r-1}$ and $v_r$ are incident with $f_r$ in
the sense of \noderef{Incident}.  These vertices are distinct by
simplicity.  Since the endpoint datum in \noderef{LooplessMultigraph} is
an unordered pair, the two distinct incidence memberships exhaust it, and
therefore
\[
 G.\operatorname{ends}(f_r)=\{v_{r-1},v_r\}.
\tag{5b}
\]
In particular, (1) and (5a) show separately that every $f_r$ differs from
the uncolored identity $e_0$; this conclusion does not identify parallel
edges by their endpoint pairs.

Reserve the last colored path identity
\[
 q:=f_p
\tag{6}
\]
for closing the trunk, and form the prefix
\[
 S=(x,e_0,y=v_1,f_2,v_2,\ldots,f_{p-1},v_{p-1}).
\tag{7}
\]
We check the clauses of \noderef{TreeSequence} for (7).  Its length is
$p-1\ge2$.  Its vertices are $x=v_p,v_1,\ldots,v_{p-1}$, hence are
pairwise distinct by the simplicity check above.  Its edge identities are
$e_0,f_2,\ldots,f_{p-1}$; the $f_r$ are pairwise distinct and each differs
from $e_0$, so these too are pairwise distinct.  Its first edge is $e_0$,
whose ends are $x,y$ as already proved.  For $2\le r\le p-1$, equation
(5b) says that the appended edge $f_r$ has new end $v_r$ and earlier end
$v_{r-1}$.  Hence each edge in (7) attaches its new displayed vertex to
an earlier displayed vertex, as required.

The sequence $S$ is also Tashkinov in the precise sense of
\noderef{TashkinovTree}.  The first identity $e_0$ has no witness
obligation.  For every later identity $f_r$ in (7), (5a) gives its color.
If that color is $\alpha$, use the earlier vertex $x$ and the first
membership in (2); if it is $\beta$, use the earlier vertex $y$ and the
second membership in (2).  Both $x$ and $y$ precede the newly attached
vertex $v_r$ because $r\ge2$.  Thus every appended colored edge has the
required earlier missing-color witness.

Now apply \noderef{ScheideMaximalExtendsEveryTree} to the $k$-triple, the
fixed-pair maximal tree $T$, and the Tashkinov tree $S$.  We obtain a
Tashkinov tree $T'$ such that
\[
 V(T')=V(T)\quad\hbox{and}\quad S\preccurlyeq T'.
\tag{8}
\]
The second assertion is the ordered relation of
\noderef{TreeSequenceInitialSegment}, not merely a set inclusion: $T'$ has
length at least $p-1$, its first $p$ vertex entries are exactly
$x,v_1,\ldots,v_{p-1}$ in the order (7), and its first $p-1$ edge entries
are exactly $e_0,f_2,\ldots,f_{p-1}$.  The first assertion in (8) is
equality of the finite sets defined by
\noderef{TreeSequenceVertexSet}.

We construct the certificate required by
\noderef{NormalGlobalMaximumTriple} and verify its fields one at a time.
Take its height to be $p$, its colors to be $\alpha,\beta$, its closing
edge to be $q$, its last index to be the index of value $p-1$ in $T'$, and
its closing path to be $P^{\mathrm{rev}}$.

First, $2\le p$ was proved above.  From $S\preccurlyeq T'$ and
$\operatorname{length}(S)=p-1$ we get
$p\le\operatorname{length}(T')+1$, so the proposed last index exists and
has value plus one equal to $p$.  The two first-vertex equalities supplied
by (8) transport (2) to the required assertions that $\alpha$ is missing
at $T'.\operatorname{vertex}(0)$ and $\beta$ at
$T'.\operatorname{vertex}(1)$.  By (5), (6), and the vertex equalities in
(8),
\[
 G.\operatorname{ends}(q)
 =\{v_{p-1},x\}
 =\{T'.\operatorname{vertex}(p-1),T'.\operatorname{vertex}(0)\},
\]
which is the closing-edge endpoint field.

Second, reversing a simple path preserves simplicity, support, and length,
so $P^{\mathrm{rev}}$ is a simple path from $y$ to $x$ of length $L$.
Since $2\mid L$ and $p=L/2+1$, its length is exactly $2(p-1)$.  Its
vertex-node support is exactly $\{v_1,\ldots,v_p\}$.  By (8), the first
$p$ vertices of $T'$ are the same identities in the reordered list
\[
 x=v_p,v_1,v_2,\ldots,v_{p-1}.
\]
Therefore, for every vertex identity $v$, the node $\mathsf{inl}(v)$ is
in the closing-path support if and only if $v$ is one of the first $p$
vertex entries of $T'$.  This proves both directions of the exact
vertex-node-support field, not only containment.

Third, the edge-node support of the reversed path is exactly
$\{f_2,\ldots,f_p\}$.  Its last member is $q=f_p$.  The remaining members
$f_2,\ldots,f_{p-1}$ are, by the edge-entry equalities in (8), exactly the
edges $T'.\operatorname{edge}(j)$ with
$1\le j$ and $j+2\le p$.  Conversely, every edge of $T'$ in that index
range is one of those displayed identities.  Hence, for every edge
identity $f$, its edge-node lies in the closing-path support if and only if
either $f=q$ or it is one of precisely those indexed trunk edges.  This is
both directions of the exact edge-node-support field.

Fourth, fix an index $j$ with $1\le j$ and $j+2\le p$.  The preserved
prefix identifies $T'.\operatorname{edge}(j)$ with $f_{j+1}$,
$T'.\operatorname{vertex}(j)$ with $v_j$, and
$T'.\operatorname{vertex}(j+1)$ with $v_{j+1}$.  Equation (5) therefore
gives
\[
 G.\operatorname{ends}(T'.\operatorname{edge}(j))
 =\{T'.\operatorname{vertex}(j),T'.\operatorname{vertex}(j+1)\}.
\]
This verifies every trunk-edge endpoint equality in the normal
certificate.

Finally, we rebuild rather than assume global maximality of $T'$.  The
$k$-triple clause is unchanged from \noderef{GlobalMaximumTriple}; the
Tashkinov-tree clause is supplied by the extension theorem; and (8) gives
\[
 |V(T')|=|V(T)|=t(G).
\]
These are exactly the three clauses of \noderef{GlobalMaximumTriple} for
$T'$.  All fields of \noderef{NormalGlobalMaximumTriple} have now been
verified, so they define normal data $N$ with
$N.\alpha=\alpha$, $N.\beta=\beta$, and
$N.\operatorname{height}=p=L/2+1$.  Only now may
\noderef{NormalTripleTrunk} be applied: it identifies the resulting trunk
with the first $p$ vertices and first $p-1$ edges already checked above.

Return the original orientation $P$ as the incidence path in the lemma's
conclusion, together with $T'$ and $N$.  Its simplicity and exact full
component support are the assertions in (3), (8) gives
$V(T')=V(T)$, and the three displayed fields of $N$ give all remaining
conjuncts.
\end{proof}
